\documentclass[11pt]{article}
% Shared print layout. The web's provisional design and MathJax are independent.
\usepackage[utf8]{inputenc}
\usepackage[T1]{fontenc}
\usepackage[spanish,es-nodecimaldot]{babel}
\usepackage{lmodern}
\usepackage{amsmath,amssymb}
\usepackage{graphicx}
\usepackage{geometry}
\geometry{a4paper,left=20mm,right=20mm,top=25mm,bottom=22mm,headheight=16pt,headsep=8mm}
\usepackage{fancyhdr,lastpage}
\usepackage{enumitem}
\usepackage{xcolor}
\usepackage{pgfplots}
\pgfplotsset{compat=1.18}
\usepgfplotslibrary{fillbetween}
\usepackage[hidelinks,pdfencoding=auto]{hyperref}
\hypersetup{pdflang=es-ES,pdfauthor={},pdfsubject={Matemáticas Aplicadas a las Ciencias Sociales II}}
\pagestyle{fancy}
\fancyhf{}
\fancyhead[L]{\small Matemáticas Aplicadas a las Ciencias Sociales II}
\fancyhead[R]{\small 2.\textsuperscript{o} Bachillerato}
\fancyfoot[L]{\footnotesize Matrices I | Edición revisada 1.0}
\fancyfoot[R]{\footnotesize Página \thepage\ de \pageref{LastPage}}
\renewcommand{\headrulewidth}{0.4pt}
\setlength{\parindent}{0pt}
\setlength{\parskip}{0.5em}
\setlist{itemsep=0.3em,topsep=0.3em}
\newcommand{\apartado}[1]{\subsection*{#1}}
\pgfplotsset{
  revision/.style={
    width=0.94\linewidth,height=7.5cm,
    axis lines=middle,grid=major,grid style={gray!18},
    tick label style={font=\small},label style={font=\small},
    xlabel={$x$},ylabel={$y$},samples=160,
    legend style={at={(0.5,-0.16)},anchor=north,draw=none,font=\small},
    clip=true
  }
}

\hypersetup{pdftitle={Matrices I - Soluciones revisadas}}
\begin{document}
\section*{Matrices I: soluciones}
Solucionario elaborado en esta revisión a partir de F014. El ZIP original no incluía soluciones.
\section*{Ejercicio 1(a): producto, diferencia e inversa}
Las dimensiones son $A:2\times3$, $B:3\times2$ y $C:2\times2$. Por tanto $AB-C$ está definida y es cuadrada.
\[
AB=\begin{pmatrix}-2&0\\6&-1\end{pmatrix},\qquad
M=AB-C=\begin{pmatrix}-3&-2\\4&2\end{pmatrix}.
\]
Por ejemplo, $(AB)_{21}=2\cdot2+(-1)(-2)+0\cdot0=6$.
\[\det M=(-3)\cdot2-(-2)\cdot4=2\neq0.\]
Existe inversa y, por la fórmula de orden 2,
\[
\boxed{M^{-1}=\frac12\begin{pmatrix}2&2\\-4&-3\end{pmatrix}
=\begin{pmatrix}1&1\\-2&-3/2\end{pmatrix}}.
\]
\textbf{Comprobación:} $MM^{-1}=M^{-1}M=I_2$.
\section*{Ejercicio 1(b): ecuación matricial}
Multiplicamos por la inversa \textbf{a la izquierda}, respetando el orden:
\[MX=C^t\quad\Longrightarrow\quad X=M^{-1}C^t.\]
Como $C$ es simétrica, $C^t=C$, y
\[
X=\begin{pmatrix}1&1\\-2&-3/2\end{pmatrix}
\begin{pmatrix}1&2\\2&-3\end{pmatrix}
=\boxed{\begin{pmatrix}3&-1\\-5&1/2\end{pmatrix}}.
\]
\textbf{Comprobación en la ecuación original:}
\[
\begin{pmatrix}-3&-2\\4&2\end{pmatrix}
\begin{pmatrix}3&-1\\-5&1/2\end{pmatrix}
=\begin{pmatrix}1&2\\2&-3\end{pmatrix}=C^t.
\]
La solución es única porque $M$ es invertible.
\clearpage
\section*{Ejercicio 2: operaciones y dimensiones}
Se muestran las cinco respuestas para consulta. En el examen se deben elegir \textbf{solo cuatro}; la puntuación total del ejercicio es 2,5 puntos.

Aquí $A,B$ son $2\times2$, $C$ es $2\times1$ y $D$ es $1\times2$. Cada operación utiliza las matrices del ejercicio 2.
\subsection*{1. $A+BC$}
El producto $BC$ existe y es $2\times1$, pero no se puede sumar a una matriz $2\times2$. La operación completa \textbf{no está definida}.
\subsection*{2. $AC+BD^t$}
\[
AC=\begin{pmatrix}8\\-17\end{pmatrix},\qquad BD^t=\begin{pmatrix}-4\\8\end{pmatrix}.
\]
Ambos son $2\times1$, luego
\[\boxed{AC+BD^t=\begin{pmatrix}4\\-9\end{pmatrix}}.\]
\subsection*{3. $B^2+CD$}
\[
B^2=\begin{pmatrix}4&0\\-8&4\end{pmatrix},\qquad
CD=\begin{pmatrix}-6&6\\-2&2\end{pmatrix}.
\]
Los dos sumandos son $2\times2$:
\[\boxed{B^2+CD=\begin{pmatrix}-2&6\\-10&6\end{pmatrix}}.\]
\subsection*{4. $|5A|$}
Las barras indican determinante. Como $A$ es de orden 2,
\[|5A|=5^2|A|=25\bigl(3\cdot1-(-1)(-6)\bigr)=25(-3)=\boxed{-75}.\]
\subsection*{5. $|D^t|$}
$D^t$ es $2\times1$ y no es cuadrada. Su determinante \textbf{no está definido}; no debe confundirse con un determinante nulo.
\clearpage
\section*{Ejercicio 3: potencia de una matriz nilpotente}
Sea
\[
N=I_3-A=\begin{pmatrix}0&-5&-6\\0&0&-7\\0&0&0\end{pmatrix}.
\]
Calculamos sus primeras potencias:
\[
N^2=\begin{pmatrix}0&0&35\\0&0&0\\0&0&0\end{pmatrix},\qquad
N^3=\begin{pmatrix}0&0&0\\0&0&0\\0&0&0\end{pmatrix}.
\]
El único término no nulo de $N^2$ es $(N^2)_{13}=(-5)(-7)=35$. Al multiplicar una vez más por $N$, también desaparece ese término.

Como $N^3=0$ y $N^2\neq0$, $N$ es nilpotente de índice 3. Para todo entero $k\geq3$, $N^k=0$.
\[
\boxed{(I_3-A)^{2025}=\begin{pmatrix}0&0&0\\0&0&0\\0&0&0\end{pmatrix}}.
\]
Se conserva el exponente \textbf{2025} del enunciado original. No hace falta calcular 2025 productos: basta demostrar que la tercera potencia es nula.
\end{document}
